\documentclass[10pt,a4paper]{amsart}
\usepackage[margin=19mm]{geometry}
\usepackage{amsmath,amssymb,mathtools}
\usepackage[scr=rsfs,cal=euler]{mathalfa}
\usepackage{xfrac}
\usepackage[unicode,hidelinks,bookmarks=false]{hyperref}
\newcommand{\ie}{i.e.\ }
\newcommand{\cf}{cf.\ }
\newcommand{\resp}{respectively\ }
\newcommand{\Subsection}{\subsection}
\providecommand{\nobreakdash}{\nobreak\hbox{-}}
\setlength{\emergencystretch}{2em}
\multlinegap0pt
\usepackage{bm}
\usepackage{bbm}
\usepackage{dsfont}
\usepackage{xspace}
\usepackage{cancel}
\usepackage{mathtools}
\usepackage{cleveref}
\usepackage{amsthm}
\makeatletter
\@ifundefined{theorem}{\newtheorem{theorem}{Theorem}}{}
\@ifundefined{lemma}{\newtheorem{lemma}[theorem]{Lemma}}{}
\@ifundefined{theorem*}{\newtheorem*{theorem*}{Theorem}}{}
\makeatother
\title[Ansatz-free Hamiltonian learning with Heisenberg-limited sca]{Ansatz-free Hamiltonian learning with Heisenberg-limited scaling}
\author{Hong-Ye Hu, Muzhou Ma, Weiyuan Gong, Qi Ye, Yu Tong, Steven T. Flammia, Susanne F. Yelin}
\date{Source: arXiv:2502.11900v2 (CC BY 4.0)}
\begin{document}
\maketitle

\noindent\emph{Source and licence.} Ansatz-free Hamiltonian learning with Heisenberg-limited scaling. Version arXiv:2502.11900v2 (CC BY 4.0), distributed under the Creative Commons Attribution 4.0 International licence, as recorded in the arXiv licence metadata for this exact version and shown on the abstract page https://arxiv.org/abs/2502.11900v2. Full rights notice: licenses/arxiv-2502.11900-CC-BY-4.0.txt.\par\smallskip

\noindent\textit{Editorial scope.} Counted unit: the Theorem* whose source label is \texttt{eq:rv\_correctness\_condition\_prob} in the article (source0.tex lines 1420--1445, source label \texttt{eq:rv\_correctness\_condition\_prob}), delivered together with the article's own complete original proof of it (source0.tex lines 1447--1492, one \texttt{proof} environment of 46 lines). No mathematics is written, repaired or re-derived: statement and proof are the article's own text line for line. Required context retained uncounted: lem:frequency\_est\_refine (lines 1393--1402); eq:rv\_correctness\_condition (lines 1397--1400). Every definition, display and label of those lines is retained unchanged. Omitted from this package and not redistributed: 1--1392, 1401--1419, 1446--1446, 1493--1930. Nothing inside a retained excerpt is omitted or altered: every retained body is delivered exactly as the article's lines are written, so no omission receipt of any kind accompanies this package. Citations: the retained excerpts carry no citation command at all, so no bibliography entry of the article is transcribed and no reference is redistributed. Reference adaptation: none was needed and none was made; every reference inside the delivered unit resolves inside it, so no unresolved reference or citation is produced. Printed slips kept literal: no printed word, symbol, hypothesis or number of the counted statement or of its proof was altered. Layout and class: the article's own class is \texttt{revtex4-2} with packages including babel, letltxmacro, latexsym, booktabs, siunitx, multirow, algpseudocode, graphicx, amsmath, amsfonts, amssymb, amsthm, cancel, mathtools; the wrapper is the lane's pinned \texttt{amsart} editorial wrapper at 10pt on A4, which restates the theorem-like environments the retained text needs and adds the article's own macro definitions that the retained text needs (none), with extra wrapper packages (bm, bbm, dsfont, xspace, cancel, mathtools, cleveref). The delivered numbering is the unit's own. Assets: no retained span contains a figure environment, a graphics-inclusion command, a PDF-page inclusion, a table or a TikZ picture; no image file is copied into this package, the package declares no asset dependency and no image file is redistributed with it. Licence: version arXiv:2502.11900v2 of this article is distributed under the Creative Commons Attribution 4.0 International licence, as recorded in the arXiv licence metadata for this exact version and shown on the abstract page https://arxiv.org/abs/2502.11900v2. A full rights notice is delivered at licenses/arxiv-2502.11900-CC-BY-4.0.txt. Characters: every retained excerpt is pure ASCII, so the delivered engine, XeLaTeX with its default Latin Modern text font, reports no missing character.
\begin{lemma}
    \label{lem:frequency_est_refine}
    Let $\theta\in[a,b]$.
    Let $Z(t)$ be a random variable such that 
    \begin{equation}
        \label{eq:rv_correctness_condition}
        |Z(t)-e^{i\theta t}|< 1/2.
    \end{equation}
    Then we can correctly distinguish between two overlapping cases $\theta\in[a,(a+2b)/3]$ and $\theta\in[(2a+b)/3,b]$ with one sample of $Z(\pi/(b-a))$. 
\end{lemma}

\begin{theorem*}
    Let $\theta\in[-A,A]$.
    Let $X(t)$ and $Y(t)$ be independent random variables satisfying
    \begin{equation}
        \label{eq:rv_correctness_condition_prob}
        \begin{aligned}
            &|X(t)-\cos(\theta t)|< 1/\sqrt{2}, \text{ with probability at least }2/3, \\
            &|Y(t)-\sin(\theta t)|< 1/\sqrt{2}, \text{ with probability at least }2/3.
        \end{aligned}
    \end{equation}
    Then with independent non-adaptive samples $X(t_1),X(t_2),\cdots,X(t_K)$ and $Y(t_1),Y(t_2),\cdots,Y(t_K)$, $t_j\geq 0$, for
    \begin{equation}
        \label{eq:number_of_samples_RPE}
        K=\mathcal{O}(\log(A/\epsilon)(\log(1/q)+\log\log(A/\epsilon))),
    \end{equation}
    \begin{equation}
    \label{eq:total_evolution_time_RPE}
    \begin{aligned}
        T=\sum_{j=1}^Kt_j=\mathcal{O}((1/\epsilon)(\log(1/q)+\log\log(A/\epsilon))),\quad \max_j t_j=\mathcal{O}(1/\epsilon),
    \end{aligned}
    \end{equation}
    we can obtain a random variable $\hat{\theta}$ such that
    \begin{equation}
        \Pr[|\hat{\theta}-\theta|>\epsilon]\leq q.
    \end{equation}
\end{theorem*}

\begin{proof}
    We let $\lambda=\epsilon/3$.
    We build a random variable $S(t)$ satisfying \eqref{eq:rv_correctness_condition}, with which we will iteratively narrow down the interval $[a,b]$ containing $\theta$ until $|b-a|\leq 2\lambda$, at which point we choose $\hat{\theta}=(a+b)/2$. If $\theta\in[a,b]$ we will then ensure $|\hat{\theta}-\theta|\leq \lambda$. However, each iteration will involve some failure probability, which we will analyze later.

    To build the random variable $S(t)$, we first use $m$ independent samples of $X(t)$ and then take median $X_{\mathrm{median}}(t)$, which satisfies 
    \[
    |X_{\mathrm{median}}(t)-\cos(\theta t)|\leq 1/\sqrt{2}
    \]
    with probability at least $1-\delta/2$, where by \eqref{eq:rv_correctness_condition_prob} and the Chernoff bound 
    \[
    \delta = c_1 e^{-c_2 m},
    \]
    for some universal constant $c_1,c_2$.
    Similarly we can obtain $Y_{\mathrm{median}}(t)$ such that
    \[
    |Y_{\mathrm{median}}(t)-\cos(\theta t)|\leq 1/\sqrt{2}
    \]
    with probability at least $1-\delta/2$. With these medians we then define
    \[
    S(t) = X_{\mathrm{median}}(t) + iY_{\mathrm{median}}(t).
    \]
    This random variable satisfies
    \[
    |S(t)-e^{i\theta t}|\leq 1/2
    \]
    with probability at least $1-\delta$ using the union bound. It therefore allows us to solve the discrimination task in \Cref{lem:frequency_est_refine} with probability at least $1-\delta$.

    Whether each iteration proceeds correctly or not, the algorithm terminates after 
    \[
    L=\lceil \log_{3/2}(A/\lambda) \rceil
    \]
    iterations. In the $l$th iteration we sample $X(s_l)$ and $Y(s_l)$ where $s_l=(\pi/2A)(3/2)^{l-1}$. 
    We use $m$ samples of $X(s_l)$ and $Y(s_l)$ for computing the median in each iteration, and therefore the failure probability is at most $\delta=c_1 e^{-c_2 m}$. 
    With probability at most $L\delta$ using the union bound, one of the iteration fails. 
    In order to ensure the protocol succeeds with probability at least $1-q$, it suffices to let $\delta \leq q/L$, and it therefore suffices to choose
    \[
    m = \lceil c_3 \log(L/q) \rceil=\mathcal{O}(\log(1/q)+\log\log(A/\epsilon)).
    \]
    The total number of samples (for either $X(t)$ or $Y(t)$) is therefore as described in \eqref{eq:number_of_samples_RPE}.
    All the $t$ in each sample added together is
    \begin{equation}
        \label{eq:total_time}
        T = m\sum_{l=1}^{L}  s_l = \frac{m\pi}{2A}\sum_{l=1}^L \left(\frac{3}{2}\right)^{l-1} \leq \frac{3\pi m}{2\epsilon}=\mathcal{O}((1/\epsilon)(\log(1/q)+\log\log(A/\epsilon))),
    \end{equation} 
    thus giving us \eqref{eq:total_evolution_time_RPE}.
\end{proof}

\end{document}
